\chapter{Context and Cover Tables}
\label{appA:chapter}

This appendix tabulates, exhaustively and with derivations, the raw
combinatorial data of the finite context site: the sixteen contexts and
their complete-block content, the generating cover families $K_A(V)$ and
$K_R(V)$ at every context, the contexts at which the two families coincide
as unmarked data, and the per-context sieve and covering-sieve counts that
the exact-topology enumeration of Chapter~11 reports. Everything here is
derivable by hand from the definitions in Chapter~10; where a value is also
registered in the executable certificate, the correspondence is noted. No
entry in this appendix carries any status beyond finite combinatorics: the
ACBE and RPM coordinate names remain formal marks with no empirical
interpretation.

\section{Naming conventions}
\label{appA:sec:naming}

\begin{notation}
\label{appA:not:names}
Atoms are $\omega_{ij} = (i, j) \in \Omega$ for $i, j \in \{0, 1\}$, where
the first coordinate is the formal ACBE mark and the second the formal RPM
mark. Blocks are $A_i = \{\omega_{i0}, \omega_{i1}\}$ (fixed first
coordinate) and $R_j = \{\omega_{0j}, \omega_{1j}\}$ (fixed second
coordinate). For the ten contexts that are neither $\varnothing$, a block,
nor $\Omega$, we use: the four singletons $\{\omega_{ij}\}$; the two
diagonals
\[
\Delta = \{\omega_{00}, \omega_{11}\}, \qquad
\nabla = \{\omega_{01}, \omega_{10}\};
\]
and the four triples $T_{ij} = \Omega \setminus \{\omega_{ij}\}$, each
named by the atom it omits. A block $B$ is \emph{complete in} $V$ when
$B \subseteq V$; the witness presheaves of Appendix~C are built on
complete-block content, so the containment columns below feed directly into
those tables.
\end{notation}

The containment rules used throughout are immediate from the definitions:
$A_0 \subseteq V$ iff $\omega_{00}, \omega_{01} \in V$;
$A_1 \subseteq V$ iff $\omega_{10}, \omega_{11} \in V$;
$R_0 \subseteq V$ iff $\omega_{00}, \omega_{10} \in V$;
$R_1 \subseteq V$ iff $\omega_{01}, \omega_{11} \in V$. Note that a
two-element context contains a complete $A$-block only if it \emph{is} that
block, and likewise for $R$-blocks; the diagonals $\Delta$ and $\nabla$
contain no complete block of either partition, because each meets each
block in exactly one atom. A triple $T_{ij}$ omits one atom and therefore
breaks exactly one $A$-block ($A_i$) and exactly one $R$-block ($R_j$),
containing the other one of each.

\section{The sixteen contexts and their complete blocks}
\label{appA:sec:contexts}

\begin{longtable}{@{}llccll@{}}
\toprule
Context $V$ & Elements & $|V|$ & Arrows into $V$ & Complete $A$-blocks &
Complete $R$-blocks \\
\midrule
\endhead
\bottomrule
\endfoot
$\varnothing$ & --- & 0 & 1 & none & none \\
$\{\omega_{00}\}$ & $\omega_{00}$ & 1 & 2 & none & none \\
$\{\omega_{01}\}$ & $\omega_{01}$ & 1 & 2 & none & none \\
$\{\omega_{10}\}$ & $\omega_{10}$ & 1 & 2 & none & none \\
$\{\omega_{11}\}$ & $\omega_{11}$ & 1 & 2 & none & none \\
$A_0$ & $\omega_{00}, \omega_{01}$ & 2 & 4 & $A_0$ & none \\
$A_1$ & $\omega_{10}, \omega_{11}$ & 2 & 4 & $A_1$ & none \\
$R_0$ & $\omega_{00}, \omega_{10}$ & 2 & 4 & none & $R_0$ \\
$R_1$ & $\omega_{01}, \omega_{11}$ & 2 & 4 & none & $R_1$ \\
$\Delta$ & $\omega_{00}, \omega_{11}$ & 2 & 4 & none & none \\
$\nabla$ & $\omega_{01}, \omega_{10}$ & 2 & 4 & none & none \\
$T_{00}$ & $\omega_{01}, \omega_{10}, \omega_{11}$ & 3 & 8 & $A_1$ & $R_1$ \\
$T_{01}$ & $\omega_{00}, \omega_{10}, \omega_{11}$ & 3 & 8 & $A_1$ & $R_0$ \\
$T_{10}$ & $\omega_{00}, \omega_{01}, \omega_{11}$ & 3 & 8 & $A_0$ & $R_1$ \\
$T_{11}$ & $\omega_{00}, \omega_{01}, \omega_{10}$ & 3 & 8 & $A_0$ & $R_0$ \\
$\Omega$ & all four & 4 & 16 & $A_0, A_1$ & $R_0, R_1$ \\
\end{longtable}

The ``arrows into $V$'' column counts the inclusions $U \subseteq V$,
i.e.\ $2^{|V|}$; summing the column gives
$1 + 4 \cdot 2 + 6 \cdot 4 + 4 \cdot 8 + 16 = 81$, the registered arrow
count of the site (Chapter~11). The complete-block columns support a
consistency count on the distribution of complete $A$-blocks over
contexts: $k = 2$ occurs only at $\Omega$ (one context); $k = 1$ occurs at
$A_0$, $A_1$, and the four triples (six contexts); $k = 0$ occurs at the
remaining nine contexts ($\varnothing$, four singletons, $R_0$, $R_1$,
$\Delta$, $\nabla$). Totals: $1 + 6 + 9 = 16$. The same distribution holds
for $R$-blocks by the evident symmetry exchanging the two coordinates.
These counts give the value cardinalities $|F_A(V)| = 2^{k}$ tabulated in
Appendix~C: nine contexts with one section, six with two, one with four.

\section{The generating family $K_A(V)$ at every context}
\label{appA:sec:ka}

$K_A(V) = \{ V \cap A_0 \to V,\; V \cap A_1 \to V \}$ at nonempty $V$,
retaining empty pieces; the null family at $\varnothing$. Each row lists
the two pieces; ``$= V$'' flags a piece equal to the whole context (an
identity arrow) and ``$= A_i$'' flags a piece that is a full block.

\begin{longtable}{@{}llll@{}}
\toprule
Context $V$ & $A_0$-piece $V \cap A_0$ & $A_1$-piece $V \cap A_1$ &
Remarks \\
\midrule
\endhead
\bottomrule
\endfoot
$\varnothing$ & \multicolumn{2}{l}{null family (no members)} &
covers by convention \\
$\{\omega_{00}\}$ & $\{\omega_{00}\} = V$ & $\varnothing$ &
identity piece present \\
$\{\omega_{01}\}$ & $\{\omega_{01}\} = V$ & $\varnothing$ &
identity piece present \\
$\{\omega_{10}\}$ & $\varnothing$ & $\{\omega_{10}\} = V$ &
identity piece present \\
$\{\omega_{11}\}$ & $\varnothing$ & $\{\omega_{11}\} = V$ &
identity piece present \\
$A_0$ & $A_0 = V$ & $\varnothing$ & identity piece present \\
$A_1$ & $\varnothing$ & $A_1 = V$ & identity piece present \\
$R_0$ & $\{\omega_{00}\}$ & $\{\omega_{10}\}$ & both proper singletons \\
$R_1$ & $\{\omega_{01}\}$ & $\{\omega_{11}\}$ & both proper singletons \\
$\Delta$ & $\{\omega_{00}\}$ & $\{\omega_{11}\}$ & both proper singletons \\
$\nabla$ & $\{\omega_{01}\}$ & $\{\omega_{10}\}$ & both proper singletons \\
$T_{00}$ & $\{\omega_{01}\}$ & $\{\omega_{10}, \omega_{11}\} = A_1$ &
both proper \\
$T_{01}$ & $\{\omega_{00}\}$ & $\{\omega_{10}, \omega_{11}\} = A_1$ &
both proper \\
$T_{10}$ & $\{\omega_{00}, \omega_{01}\} = A_0$ & $\{\omega_{11}\}$ &
both proper \\
$T_{11}$ & $\{\omega_{00}, \omega_{01}\} = A_0$ & $\{\omega_{10}\}$ &
both proper \\
$\Omega$ & $A_0$ & $A_1$ & both proper: the global cover \\
\end{longtable}

Two structural observations, both used in the body of the monograph.
First, the generated sieve is nonmaximal exactly when \emph{no} piece is
an identity: this happens at $R_0$, $R_1$, $\Delta$, $\nabla$, the four
triples, and $\Omega$ (ten contexts). At the four singletons and at $A_0$,
$A_1$, one piece is the identity, the generated sieve is maximal, and the
cover imposes no sheaf constraint beyond functoriality; at $\varnothing$
the null cover imposes the one-section constraint. The certificate check
\texttt{proper\_global\_covers} registers the $\Omega$ row. Second, at
every context each nonempty piece lies inside a single block, so the
singleton-factor bijection of Chapter~10 applies to it; this is visible in
the table as the absence of any piece meeting both blocks.

\section{The generating family $K_R(V)$ at every context}
\label{appA:sec:kr}

$K_R(V) = \{ V \cap R_0 \to V,\; V \cap R_1 \to V \}$, same conventions.

\begin{longtable}{@{}llll@{}}
\toprule
Context $V$ & $R_0$-piece $V \cap R_0$ & $R_1$-piece $V \cap R_1$ &
Remarks \\
\midrule
\endhead
\bottomrule
\endfoot
$\varnothing$ & \multicolumn{2}{l}{null family (no members)} &
covers by convention \\
$\{\omega_{00}\}$ & $\{\omega_{00}\} = V$ & $\varnothing$ &
identity piece present \\
$\{\omega_{01}\}$ & $\varnothing$ & $\{\omega_{01}\} = V$ &
identity piece present \\
$\{\omega_{10}\}$ & $\{\omega_{10}\} = V$ & $\varnothing$ &
identity piece present \\
$\{\omega_{11}\}$ & $\varnothing$ & $\{\omega_{11}\} = V$ &
identity piece present \\
$A_0$ & $\{\omega_{00}\}$ & $\{\omega_{01}\}$ & both proper singletons \\
$A_1$ & $\{\omega_{10}\}$ & $\{\omega_{11}\}$ & both proper singletons \\
$R_0$ & $R_0 = V$ & $\varnothing$ & identity piece present \\
$R_1$ & $\varnothing$ & $R_1 = V$ & identity piece present \\
$\Delta$ & $\{\omega_{00}\}$ & $\{\omega_{11}\}$ & both proper singletons \\
$\nabla$ & $\{\omega_{10}\}$ & $\{\omega_{01}\}$ & both proper singletons \\
$T_{00}$ & $\{\omega_{10}\}$ & $\{\omega_{01}, \omega_{11}\} = R_1$ &
both proper \\
$T_{01}$ & $\{\omega_{00}, \omega_{10}\} = R_0$ & $\{\omega_{11}\}$ &
both proper \\
$T_{10}$ & $\{\omega_{00}\}$ & $\{\omega_{01}, \omega_{11}\} = R_1$ &
both proper \\
$T_{11}$ & $\{\omega_{00}, \omega_{10}\} = R_0$ & $\{\omega_{01}\}$ &
both proper \\
$\Omega$ & $R_0$ & $R_1$ & both proper: the global cover \\
\end{longtable}

The table is the image of the $K_A$ table under the coordinate exchange
$(i, j) \mapsto (j, i)$, which swaps $A$-blocks with $R$-blocks, fixes
$\varnothing$, $\Delta$, and $\Omega$, and permutes the remaining contexts
accordingly ($A_i \leftrightarrow R_i$, $\nabla$ fixed as a set with its
atoms swapped, $T_{01} \leftrightarrow T_{10}$). This symmetry is the
tabular form of the fact, used in Chapter~10, that every ACBE statement
has an exact RPM mirror.

\section{Coincidence of the unmarked families}
\label{appA:sec:coincide}

The marks chapter and the certificate checks
\texttt{singleton\_unmarked\_families\_coincide} /
\texttt{singleton\_marked\_covers\_distinct} turn on the question: at
which contexts do $K_A(V)$ and $K_R(V)$ carry the same information once
the block labels are stripped? Comparing the two tables row by row, as
\emph{unordered sets of member arrows} (the exact list the topology
consumes, forgetting which block indexed which piece):

\begin{longtable}{@{}lllc@{}}
\toprule
Context $V$ & $K_A(V)$ as arrow set & $K_R(V)$ as arrow set &
Coincide? \\
\midrule
\endhead
\bottomrule
\endfoot
$\varnothing$ & null & null & yes \\
$\{\omega_{ij}\}$ (each of 4) & $\{V, \varnothing\}$ &
$\{V, \varnothing\}$ & yes \\
$A_0$ & $\{A_0, \varnothing\}$ &
$\{\{\omega_{00}\}, \{\omega_{01}\}\}$ & no \\
$A_1$ & $\{A_1, \varnothing\}$ &
$\{\{\omega_{10}\}, \{\omega_{11}\}\}$ & no \\
$R_0$ & $\{\{\omega_{00}\}, \{\omega_{10}\}\}$ &
$\{R_0, \varnothing\}$ & no \\
$R_1$ & $\{\{\omega_{01}\}, \{\omega_{11}\}\}$ &
$\{R_1, \varnothing\}$ & no \\
$\Delta$ & $\{\{\omega_{00}\}, \{\omega_{11}\}\}$ &
$\{\{\omega_{00}\}, \{\omega_{11}\}\}$ & yes \\
$\nabla$ & $\{\{\omega_{01}\}, \{\omega_{10}\}\}$ &
$\{\{\omega_{01}\}, \{\omega_{10}\}\}$ & yes \\
$T_{00}$ & $\{\{\omega_{01}\}, A_1\}$ & $\{\{\omega_{10}\}, R_1\}$ & no \\
$T_{01}$ & $\{\{\omega_{00}\}, A_1\}$ & $\{R_0, \{\omega_{11}\}\}$ & no \\
$T_{10}$ & $\{A_0, \{\omega_{11}\}\}$ & $\{\{\omega_{00}\}, R_1\}$ & no \\
$T_{11}$ & $\{A_0, \{\omega_{10}\}\}$ & $\{R_0, \{\omega_{01}\}\}$ & no \\
$\Omega$ & $\{A_0, A_1\}$ & $\{R_0, R_1\}$ & no \\
\end{longtable}

\begin{finding}
\label{appA:fnd:coincide}
The unmarked families coincide at exactly seven of the sixteen contexts:
the empty context, the four singletons, and the two diagonals $\Delta$ and
$\nabla$. At these contexts the generated sieves are equal and the
topology retains no trace of which partition produced the cover; only the
marked records (different cover kinds, different digests) distinguish
them. At the remaining nine contexts the arrow sets differ and the
topology itself separates the covers. Three refinements of this statement
are worth recording. (i) The singleton and empty cases are the ones
registered as certificate checks; the diagonal cases are derived here
directly from the tables and are not separately registered in v0.1.
(ii) At $\Delta$ the coincidence is aligned block-by-block: the
$A_0$-piece equals the $R_0$-piece ($\{\omega_{00}\}$) and the $A_1$-piece
equals the $R_1$-piece ($\{\omega_{11}\}$). At $\nabla$ the coincidence
is \emph{crossed}: the $A_0$-piece equals the $R_1$-piece
($\{\omega_{01}\}$) and the $A_1$-piece equals the $R_0$-piece
($\{\omega_{10}\}$). So even among coinciding contexts, the
indexed-by-block families agree at $\Delta$-type contexts only up to
relabeling and at $\nabla$ only up to a label crossing --- a second,
finer layer of information the unmarked arrow set forgets. (iii) The
coincidence pattern is forced by block geometry: $K_A(V) = K_R(V)$ as
arrow sets can happen only when every piece of both covers has at most
one atom, i.e.\ when $V$ meets each block in at most one atom, which
characterizes exactly the seven listed contexts (the subsets of $\Omega$
containing no complete block and no two atoms sharing a block occur
precisely at sizes $\le 1$ and at the two diagonals).
\end{finding}

\section{The common refinement at every context}
\label{appA:sec:refinement}

The common refinement of $K_A(V)$ and $K_R(V)$ has as members the four
pairwise intersections $V \cap A_i \cap R_j$, $(i, j) \in \{0,1\}^2$.
By transversality, $V \cap A_i \cap R_j = V \cap \{\omega_{ij}\}$: the
member indexed by $(i, j)$ is $\{\omega_{ij}\}$ when
$\omega_{ij} \in V$ and $\varnothing$ otherwise. The refinement is thus
the singleton family of $V$ padded with empty pieces --- the
context-by-context content of the certificate check
\texttt{transverse\_refinement\_is\_singletons} and, at $\Omega$, of
\texttt{joint\_refinement\_is\_four\_singletons}.

\begin{longtable}{@{}llc@{}}
\toprule
Context $V$ & Nonempty refinement members & Empty members \\
\midrule
\endhead
\bottomrule
\endfoot
$\varnothing$ & none & 4 \\
$\{\omega_{ij}\}$ (each of 4) & $\{\omega_{ij}\}$ & 3 \\
$A_0$ & $\{\omega_{00}\}, \{\omega_{01}\}$ & 2 \\
$A_1$ & $\{\omega_{10}\}, \{\omega_{11}\}$ & 2 \\
$R_0$ & $\{\omega_{00}\}, \{\omega_{10}\}$ & 2 \\
$R_1$ & $\{\omega_{01}\}, \{\omega_{11}\}$ & 2 \\
$\Delta$ & $\{\omega_{00}\}, \{\omega_{11}\}$ & 2 \\
$\nabla$ & $\{\omega_{01}\}, \{\omega_{10}\}$ & 2 \\
$T_{ij}$ (each of 4) & the three singletons of $T_{ij}$ & 1 \\
$\Omega$ & all four singletons & 0 \\
\end{longtable}

Empty members and their repetitions do not change the generated sieve,
so at every context the refinement generates the sieve of size-$\le 1$
subcontexts --- the generating sieve of $J_{\mathrm{atom}}$. This is
the per-context form of the join computation in Chapter~10, and it also
explains the coincidence pattern of
Finding~\ref{appA:fnd:coincide} from another angle: at the seven
coinciding contexts, $K_A(V)$ and $K_R(V)$ each \emph{already equal}
their own common refinement up to empty members, so nothing remains to
distinguish them at the arrow level.

\section{Per-context sieve and covering-sieve counts}
\label{appA:sec:sievecounts}

For completeness we tabulate the sieve universe and the joined-topology
covering counts underlying the registered totals $298$ and $168$
(Chapter~11). A sieve on $V$ is a down-set of $\mathcal{P}(V)$, counted by
the Dedekind number $M(|V|)$; a sieve is covering for
$J_A \vee J_R = J_{\mathrm{atom}}$ iff it contains every singleton of $V$
(Chapter~10).

\begin{longtable}{@{}lcccc@{}}
\toprule
Context class & Count of contexts & Sieves on each & Covering (join) on
each & Subtotals (sieves / covering) \\
\midrule
\endhead
\bottomrule
\endfoot
$|V| = 0$ & 1 & $M(0) = 2$ & 2 & 2 / 2 \\
$|V| = 1$ & 4 & $M(1) = 3$ & 1 & 12 / 4 \\
$|V| = 2$ & 6 & $M(2) = 6$ & 2 & 36 / 12 \\
$|V| = 3$ & 4 & $M(3) = 20$ & 9 & 80 / 36 \\
$|V| = 4$ & 1 & $M(4) = 168$ & 114 & 168 / 114 \\
\midrule
Totals & 16 & --- & --- & 298 / 168 \\
\end{longtable}

Derivation notes for the covering column. At $\varnothing$ both sieves
cover: the null cover generates the empty sieve, and its unique supersieve
is maximal. At a singleton the only covering sieve is maximal (the sole
generating family contains the identity piece). At $|V| = 2$ the covering
sieves are the sieve of all proper subcontexts and the maximal sieve. At
$|V| = 3$, a covering sieve contains all four subcontexts of size
$\le 1$; the free data is a down-set over the three $2$-subsets and the
top, giving $2^3 = 8$ choices without the top plus $1$ with it (the top
forces all $2$-subsets), total $9$. At $\Omega$, the free data is a
down-set of the eleven subcontexts of size $\ge 2$; counting by the set of
included $3$-subsets (each forcing its three $2$-subsets, with any two
$3$-subsets sharing one $2$-subset and any three covering all six) gives
$64 + 32 + 12 + 4 + 1 = 113$ down-sets without $\Omega$ itself plus $1$
with it, total $114$. These hand counts reproduce the enumeration totals
that the modality-composition verifier computes independently and reports
identically for both presentations of the joined topology.
